```latex
\documentclass[12pt]{article}
\usepackage[utf8]{inputenc}
\usepackage{amsmath, amssymb, graphicx, hyperref}
\usepackage{geometry}
\geometry{a4paper, margin=1in}
\title{SHOA-Praktik: Neuro Chip DevKit -- A Neuromorphic Hardware Platform with Intrinsic Fault Learning}
\author{Konstantin Fedotov}
\date{\today}

\begin{document}

\maketitle

\begin{abstract}
We present the SHOA-Praktik Neuro Chip DevKit, the fifth practical module of the SHOA research program. Building on the SHOA-Neuro (DOI: 10.5281/zenodo.20033465) and NeuroSHOA (DOI: 10.5281/zenodo.20267413) preprints, the DevKit provides a complete hardware and software platform for developing neuromorphic chips that ``learn from errors''. The core innovation is a memristor array integrated with CdSe/ZnS quantum dots that stores ``pre-failure memory'' and is activated by a Grape pulse. The kit includes an FPGA-based DevBoard, an SDK with a fault-learning emulator, and a calibration suite. We provide four formal validation criteria: recovery time, return accuracy, fault resilience over cycles, and energy overhead. The Neuro Chip DevKit transforms the SHOA theory of spin-based learning into a silicon-ready, testable engineering platform for resilient AI hardware.
\end{abstract}

\section{Introduction}
Modern AI accelerators and neuromorphic chips face a critical bottleneck: they are brittle. A single hardware fault, a radiation-induced bit flip, or a timing error can lead to silent data corruption or catastrophic system failure. Traditional fault tolerance relies on static redundancy (e.g., triple modular redundancy), which incurs high area and energy costs.

The SHOA project has established an alternative paradigm: dynamic, intrinsic learning from faults. The SHOA-Neuro model proved that an ensemble of quantum dots with spin memory can adaptively lower its voting threshold after each recovery cycle, mimicking biological long-term potentiation (LTP). The Neuro Chip DevKit is the physical embodiment of this principle, providing a complete platform for building and testing chips that become more resilient over time.

\section{Design and Architecture}
\subsection{System Components}
The Neuro Chip DevKit consists of five integrated subsystems:
\begin{itemize}
    \item \textbf{Quantum Dot Neuromorphic Elements:} Memristors embedded with CdSe/ZnS quantum dots (4--6 nm). These elements serve dual functions: storing ``pre-failure memory'' and acting as resonators for the Grape pulse. The target density is up to $10^9$ elements per cm$^2$.
    \item \textbf{Grape Pulse Interface:} A dual-mode communication channel. The optical mode uses a femtosecond laser pulse (100--500 fs) for fast, global activation. The electrical mode uses spike-based signals between neurons. Both operate under the SHOA-Neuro Pulse Protocol (SNPP).
    \item \textbf{Fault Learning Emulator:} A software environment that simulates hardware faults (overheating, radiation, data errors) and tracks the chip's adaptive response over multiple cycles. It provides real-time visualization of learning curves.
    \item \textbf{DevBoard and SDK:} The hardware platform is based on an FPGA (Xilinx Zynq) or ARM Cortex-M7, with USB 3.0, Ethernet, and HDMI interfaces. The SDK supports Python, C++, and Verilog, and includes libraries for Spiked Voting (SHOA-Neuro Core), pulse generation (Grape-Pulse Engine), and fault simulation (Fault-Learning Toolkit).
    \item \textbf{Calibration and Testing Software:} A Python/C++ application for tuning Grape pulse parameters, calibrating memristor states, and running automated stress tests with metric collection.
\end{itemize}

\subsection{Operational Workflow}
The fault-learning cycle follows a strict sequence:
\begin{enumerate}
    \item \textbf{Initialization:} The chip is loaded with a reference state (``pre-failure memory'').
    \item \textbf{Fault Injection:} The emulator or a real environmental stressor introduces a fault (e.g., overheated node, corrupted weight).
    \item \textbf{Detection and Isolation:} The monitoring system detects the deviation. The faulty node is isolated by gating its output.
    \item \textbf{Grape Pulse and Voting:} A Grape pulse (optical or electrical) is applied. The Spiked Voting protocol selects the most stable state from the pre-failure memory, and the chip's weights are restored.
    \item \textbf{Learning:} The chip adjusts its effective voting threshold $K_{\text{eff}}$ according to the memory law $K_{\text{eff}}(N) = K_0(1 - \beta e^{-\lambda N_{\text{recovery}}})$, becoming faster and more efficient with each cycle.
\end{enumerate}

\section{Formal Validation Criteria}
To ensure that the DevKit demonstrates genuine SHOA-Neuro fault learning, we impose the following mandatory criteria.

\subsection{Criterion I: Recovery Time}
\begin{itemize}
    \item \textbf{Definition:} The chip must recover from a single-node hardware fault in under 1 ms.
    \item \textbf{Measurement:} Inject a stuck-at fault in a memristor and measure the time from detection to full restoration of correct computation.
    \item \textbf{Target:} $t_{\text{recovery}} < 1$ ms.
\end{itemize}

\subsection{Criterion II: Return Accuracy}
\begin{itemize}
    \item \textbf{Definition:} After recovery, the chip's computational accuracy must return to over 95\% of its pre-fault baseline.
    \item \textbf{Measurement:} Run a standard MNIST or CIFAR-10 classification task before the fault and after recovery.
    \item \textbf{Target:} Accuracy$_{\text{post}} > 0.95 \times$ Accuracy$_{\text{pre}}$.
\end{itemize}

\subsection{Criterion III: Cyclic Fault Resilience}
\begin{itemize}
    \item \textbf{Definition:} The chip must maintain over 99\% correct operation after 100 consecutive fault-recovery cycles.
    \item \textbf{Measurement:} Run an automated loop that injects and recovers from a fault 100 times, tracking the number of operations that produce correct results.
    \item \textbf{Target:} 99\% of operations correct after 100 cycles.
\end{itemize}

\subsection{Criterion IV: Energy Overhead}
\begin{itemize}
    \item \textbf{Definition:} The Grape pulse and voting mechanism must consume less than 5\% additional energy compared to the chip's baseline operational power.
    \item \textbf{Measurement:} Measure the total energy consumption of the DevBoard over a 1-hour period with and without the SHOA recovery mechanisms active (but without injected faults).
    \item \textbf{Target:} $P_{\text{SHOA}} < 1.05 \times P_{\text{baseline}}$.
\end{itemize}

\section{Conclusion}
The SHOA-Praktik Neuro Chip DevKit is the first hardware platform to demonstrate intrinsic fault learning in a neuromorphic architecture. By satisfying the four formal criteria, any research lab can independently verify that their neuromorphic design achieves SHOA-compliant resilience, paving the way for a new generation of self-healing AI hardware.

\end{document}
